\begin{abstract}
Although large unsupervised language models (LMs) demonstrate strong capabilities in acquiring extensive world knowledge and reasoning abilities, precisely shaping their outputs remains challenging because their training is entirely unsupervised. To introduce greater controllability, current approaches gather human judgments comparing model outputs and subsequently fine-tune the LM to better match these preferences, frequently using reinforcement learning from human feedback (RLHF). Nevertheless, RLHF involves a multi-step and often delicate process: it first builds a reward model capturing human preferences, and then applies reinforcement learning to adapt the LM to optimize this inferred reward, all while trying not to deviate excessively from the pretrained model. In this study, we present a novel formulation of the reward model within RLHF that makes it possible to derive the associated optimal policy analytically, thus permitting the RLHF objective to be addressed via a straightforward classification loss. This new technique, termed \textit{Direct Preference Optimization} (DPO), is robust, effective, and resource-efficient—removing the need for sampling from the LM during tuning and reducing hyperparameter tuning overhead. Experimental results indicate that DPO fine-tunes LMs to follow human preferences at least as well as, and in some cases better than, leading existing strategies. Specifically, DPO not only surpasses PPO-based RLHF in modulating the sentiment of generated text, but also matches or surpasses its performance in tasks such as summarization and single-turn dialogue, all while offering greater simplicity in implementation and training.
\end{abstract}

\section{Introduction}
Recent advancements in unsupervised language models (LMs) trained on massive corpora have revealed unexpected proficiencies~\citep{chowdhery2022palm, brown2020language, touvron2023llama,bubeck2023sparks}. Nevertheless, these LMs are exposed to text authored by people with diverse intentions, values, and expertise. Certain tendencies found in the training data may be inappropriate as behavioral templates; for instance, it is beneficial if an AI code assistant can \textit{recognize} typical programming errors to address them, but when actually writing code, we wish to shift its output distribution in favor of the (perhaps infrequent) instances of exemplary programming included in the corpus. Analogously, we expect a language model to be \textit{informed} of prevalent misconceptions—such as a fallacy believed by 50\% of individuals—yet, we do not desire the model to propagate this misinformation in half its related answers. Thus, ensuring that the model's \emph{output aligns with preferred conduct and content}, as opposed to simply mirroring its vast \textit{acquired knowledge and skills}, is fundamental for developing AI systems that are reliable, effective, and controllable \citep{ouyang2022training}. Traditional techniques for steering LMs toward desirable outcomes utilize reinforcement learning (RL) with human feedback, but we demonstrate that the optimization target adopted in these RL approaches can be reformulated exactly as a simple binary cross-entropy loss, which leads to a substantial simplification of the preference alignment workflow.

Conceptually, current frameworks promote favorable behaviors in a language model by leveraging curated datasets capturing human preferences for traits such as safety and usefulness. This preference alignment phase succeeds an initial stage involving broad, unsupervised pre-training over extensive textual material. While preference learning can be approached via direct supervised fine-tuning on high-quality human-generated answers, reinforcement learning from human or artificial feedback (RLHF/RLAIF; \citep{christiano2017deep,bai2022constitutional}) has achieved the most success. RLHF protocols entail first training a reward model from annotated human preference data and subsequently employing RL to fine-tune a language model so that its policy produces responses scored highly by the reward model, while maintaining proximity to the original LM's behavior. Although RLHF yields models with remarkable conversational and programming skills, this approach is notably more intricate than conventional supervised training, as it necessitates training several language models and sampling from their policies during optimization, thereby imposing considerable computational overhead.

In this work, we introduce a method for optimizing language models to align with human preferences directly, circumventing the need for explicit reward modeling or reinforcement learning approaches. We present \textit{{Direct Preference Optimization} (DPO)}, an algorithm that inherently optimizes the same target as conventional RLHF approaches (maximization of reward under a KL-divergence regularization), yet remains both easy to implement and efficient to train. At a high level, DPO modifies the log-probability ratio between responses favored and not favored according to preferences, but crucially utilizes an adaptive, per-example weighting mechanism. This adjustment guards against the model collapse that can emerge when employing a simplistic probability ratio formulation. As with other methods, DPO leverages a formal preference model—such as the Bradley-Terry model \cite{bradley1952rankanalysis}—to evaluate the alignment of a proposed reward function with observed preference labels. Notably, traditional algorithms employ this model to craft a preference-based loss for reward model training before proceeding to train a policy on this learned reward. In contrast, DPO capitalizes on a variable transformation to express the preference loss directly in terms of the policy. This allows, given human-labeled preference datasets over model outputs, the policy to be updated via a binary cross entropy loss, efficiently learning the policy corresponding to an implicit reward model that best matches the preference data.

The principal innovation we put forth is Direct Preference Optimization (DPO): a straightforward, reinforcement learning-free approach for preference-based training of language models. Experimental results indicate that DPO performs on par with state-of-the-art methods, including those based on PPO-driven RLHF, across several preference learning benchmarks such as sentiment steering, text summarization, and conversational modeling, utilizing language models scaling up to 6B parameters.

\section{Related Work}

Self-supervised language models at larger scales have demonstrated the capacity to solve various tasks without explicit training examples, utilizing either zero-shot \citep{radford2019language} or few-shot prompting techniques \citep{gpt3,megatron,chowdhery2022palm}. Nevertheless, their effectiveness on specific applications and their alignment with user intentions can be markedly enhanced through fine-tuning on curated datasets containing explicit instructions paired with human-authored responses \citep{mishra-etal-2022-cross,sanh2022multitask,chung2022scaling,thoppilan2022lamda}. This process, known as `instruction-tuning', allows large language models (LLMs) to extrapolate to novel instructions beyond those seen during tuning and improves overall user experience \citep{chung2022scaling}. Even with the successes of instruction-tuning, obtaining \textit{relative} human assessments of response quality is generally less demanding than sourcing expert-level demonstrations. As a result, more recent approaches have concentrated on fine-tuning LLMs using collections of human preference data, leading to advancements in areas such as translation \citep{kreutzer-etal-2018-reliability}, summarization \citep{stiennon2022learning,ziegler2020finetuning}, narrative generation \citep{ziegler2020finetuning}, and instruction compliance \citep{ouyang2022training,ramamurthy2023is}. These strategies typically involve initially learning a reward model that predicts preference alignment—using, for example, the Bradley-Terry model \citep{bradley1952rankanalysis}—and subsequently training the language model to maximize this reward by means of reinforcement learning methods such as REINFORCE \citep{williams1992reinforce}, proximal policy optimization (PPO; \cite{schulman2017proximal}), or their extensions \citep{ramamurthy2023is}. Related research further utilizes LLMs already tuned with human feedback to automatically produce more preference data, targeting specific properties like safety or innocuousness \citep{bai2022constitutional}, with minimal direct human labeling, instead relying on textual guidelines for model self-assessment. Collectively, these methods fuse advances from two key research directions: one focusing on applying reinforcement learning to language model optimization across diverse goals~\citep{Ranzato2015SequenceLT,paulus2018a,wu2018learning}, and another advancing frameworks for training from human preference feedback \citep{christiano2017deep,kupcsik2018learning}. While utilizing relative preferences from humans is highly attractive, the practicalities of reinforcement learning-based fine-tuning for LLMs introduce substantial difficulties; in response, the present work introduces a theoretically-grounded method for optimizing over relative preferences, avoiding reliance on reinforcement learning.

Beyond applications involving language, the task of inferring policies from preferences has been addressed within both bandit and reinforcement learning frameworks, leading to the development of a range of methodologies. When action preferences or rankings serve as feedback in place of explicit rewards, this formulation is referred to as the contextual dueling bandit (CDB; \cite{yue2012karmed,dudik2015contextual}). Because absolute rewards are unavailable in CDBs, theoretical analyses rely on the concept of a \textit{von Neumann winner}, defined as a policy that achieves an expected win rate of no less than 50\% when compared against \textit{any} alternative policy \citep{dudik2015contextual}. Notably, CDBs operate in an online setting where preference feedback arrives in real time, contrasting with scenarios in which human-derived preferences are provided as a static offline batch of annotated action pairs \citep{yan2022human}. Likewise, in \textit{preference-based RL} (PbRL), learning is driven by binary preferences determined by an unknown 'scoring' function rather than direct rewards \citep{BusaFekete2014,ruiz2023dueling}. While numerous PbRL algorithms are available—including those capable of leveraging off-policy preference data—the prevailing paradigm generally entails first constructing an explicit model of the latent scoring (reward) function before performing policy optimization \citep{jain2013learning,BusaFekete2014,christiano2017deep,sadigh2017active,kupcsik2018learning}. In contrast, our work proposes a unified policy learning approach that eschews reward model estimation and instead directly optimizes the policy to align with preferences.

\section{Preliminaries}\label{section:prelims}

We briefly summarize the RLHF workflow as outlined by \citeauthor{ziegler2020finetuning} and subsequently followed by others \citep{stiennon2022learning, bai2022training, ouyang2022training}. The process is typically divided into three main steps: 1) supervised fine-tuning (SFT); 2) collection of preference data and reward model training; and 3) reinforcement learning (RL)-based policy optimization.

\textbf{SFT}: The RLHF process typically initiates by applying supervised learning to a pre-trained language model using curated datasets that are closely aligned with the end-task objectives (e.g., conversational agents, summarization systems), resulting in a supervised fine-tuned model, denoted as $\pi^\text{SFT}$.

\textbf{Reward Modelling Phase}: During this stage, the SFT model is conditioned on prompts $x$ to generate pairs of candidate responses $(y_1, y_2)\sim \pi^\text{SFT}(y \mid x)$. Human annotators then review each pair, indicating a preference for one response over the other; this is denoted by $y_w\succ y_l \mid x$, where $y_w$ is the selected (preferred) response and $y_l$ the less preferred one. These pairwise preferences are presumed to originate from a hidden reward function $r^*(y, x)$, to which we do not have direct access. Several probabilistic models can represent human preferences, with the Bradley-Terry (BT) model \cite{bradley1952rankanalysis} being among the most frequently employed; more broadly, Plackett-Luce models \citep{plackett1975analysis, luce2012individual} are suitable if ranked lists of outputs are available. The BT model describes the distribution over human preferences $p^*$ as follows:

Given a fixed dataset of pairwise comparisons $\mathcal{D} = \{x^{(i)}, y_w^{(i)}, y_l^{(i)}\}_{i=1}^N$ drawn according to $p^*$, a reward function $r_\phi(x, y)$ may be parameterized, and its parameters inferred through maximizing the likelihood. When the task is interpreted as binary classification, the objective to be minimized is the negative log-likelihood loss:

\begin{equation}\label{eq:reward_model}
    \mathcal{L}_R(r_{\phi}, \mathcal{D}) = -\mathbb{E}_{(x, y_w, y_l)\sim \mathcal{D}}\left[\log \sigma(r_{\phi}(x, y_w)- r_{\phi}(x, y_l))\right]
\end{equation}

where $\sigma$ denotes the logistic sigmoid function. In language model applications, the reward network $r_{\phi}(x, y)$ commonly initializes from the SFT model $\pi^\text{SFT}(y \mid x)$, with a linear output head added on top of the terminal transformer layer to produce a scalar reward prediction \cite{ziegler2020finetuning}. To limit variance in the learned rewards, previous studies normalize rewards to ensure $\mathbb{E}_{x,y\sim \mathcal{D}}\left[r_\phi(x, y)\right] = 0$ for every $x$.

\textbf{RL Fine-Tuning Phase}: In the RL fine-tuning step, feedback from the trained reward model is provided to guide the policy of the language model. As established in prior work~\citep{jaques2017sequence, jaques2020human}, the optimization can be formulated as:

\begin{equation}\label{eq:RL}
\max_{\pi_{\theta}}  \mathbb{E}_{x\sim \mathcal{D}, y\sim \pi_{\theta}(y \mid x)}\left[r_{\phi}(x, y)\right] - \beta\mathbb{D}_{\textrm{KL}}\left[\pi_{\theta}(y\mid x)\mid \mid \pi_\text{ref}(y\mid x)\right],
\end{equation}

where $\beta$ controls the strength of the regularization that enforces closeness to the reference policy $\pi_\text{ref}$, which is typically taken to be the initial SFT policy $\pi^\text{SFT}$. Commonly, the model policy $\pi_\theta$ is initialized with parameters from $\pi^\text{SFT}$. This regularization is necessary: it constrains the learned policy from straying far from domains where the reward estimator remains reliable, and helps avoid issues such as loss of diversity or convergence to degenerate policies that yield only a few high-reward responses. Since text generation is a discrete process, the gradient of this objective cannot be computed directly, necessitating reinforcement learning for optimization. The dominant approach in the literature \citep{ziegler2020finetuning, stiennon2022learning, bai2022training, ouyang2022training} is to define the reward as ${r(x, y) = r_{\phi}(x, y) -\beta (\log \pi_{\theta}(y\mid x) - \log \pi_\text{ref}(y\mid x))}$, and employ PPO \cite{schulman2017proximal} for maximization.

\section{Direct Preference Optimization}\label{sec:DPO}

The practical difficulties of applying reinforcement learning at the scale required for language model fine-tuning motivate the development of a more streamlined method for optimizing policies with preference data. Instead of first learning a reward model and subsequently optimizing it with RL—an approach common to previous RLHF methodologies—our method adopts a parameterization of the reward such that the policy optimal with respect to the reward can be obtained analytically, thereby removing the need for iterative RL-based optimization. As discussed in detail below, this approach capitalizes on an explicit mapping from the reward function to the corresponding optimal policy, which in turn allows the transformation of an objective on reward models into an equivalent objective on policy models. By leveraging this reparameterization, we can directly optimize the policy under preference modeling frameworks (such as the Bradley-Terry model), without the intermediate step of fitting a separate reward model. Here, the policy itself both generates responses and implicitly encodes preferences via its parameters.

\textbf{Derivation of the DPO Objective.} We begin by considering the standard RL objective, Eq.~\ref{eq:RL}, framed with a generic reward function $r$, as established in existing literature. As demonstrated in earlier works~\citep{peters2007reinforcement, peng2019advantage, korbak2022reinforcement, go2023aligning}, it can be directly shown that the optimal policy for maximizing the KL-regularized reward (see Eq.~\ref{eq:RL}) assumes the following form:

\begin{equation}\label{eq:op_policy}
    \pi_r(y\mid x) = \frac{1}{Z(x)}\pi_\text{ref}(y\mid x)\exp\left(\frac{1}{\beta}r(x, y)\right),
\end{equation}

where the normalization constant $Z(x) = \sum_{y}\pi_\text{ref}(y\mid x)\exp\left(\frac{1}{\beta}r(x, y)\right)$ ensures a valid probability distribution. A detailed step-by-step derivation is given in Appendix~\ref{app:derivation1}. Nonetheless, directly evaluating $Z(x)$ remains computationally intensive, even when approximating the true reward function $r^*$ using a maximum likelihood estimate $r_\phi$~\citep{korbak2022reinforcement, go2023aligning}. This computational burden complicates the practical utility of the above form. 

To address this, we can algebraically manipulate Eq.~\ref{eq:op_policy} to re-express the reward function as a function of its corresponding optimal policy $\pi_r$, the reference distribution $\pi_\text{ref}$, and the, in general, unknown partition function $Z(\cdot)$. Specifically, by applying the logarithm to both sides of Eq.~\ref{eq:op_policy} and rearranging, we have:

\begin{equation}\label{eq:main_eq}
    r(x,y) =\beta \log \frac{\pi_r(y\mid x)}{\pi_\text{ref}(y\mid x)} + \beta \log Z(x).
\end{equation}

This transformation can equivalently be applied to the true reward $r^*$ and the optimal policy $\pi^*$. Importantly, under the Bradley-Terry framework, outcomes depend solely on differences in rewards between two completions: ${p^*(y_1 \succ y_2 \mid x) = \sigma(r^*(x, y_1) - r^*(x, y_2))}$. Plugging the parameterization from Eq.~\ref{eq:main_eq} into the preference model Eq.~\ref{eq:bradley-terry} for $r^*(x, y)$ eliminates the partition function, resulting in a preference probability expression that is solely a function of $\pi^*$ and $\pi_\text{ref}$. Thus, the policy $\pi^*$ derived from RLHF in the Bradley-Terry setting is governed by the following preference model:

\begin{equation}\label{eq:objective}
    p^*(y_1\succ y_2 \mid x)=\frac{1}{1 + \exp\left(\beta \log \frac{\pi^*(y_2\mid x)}{\pi_\text{ref}(y_2\mid x)} - \beta \log \frac{\pi^*(y_1\mid x)}{\pi_\text{ref}(y_1\mid x)}\right)}
\end{equation}

A more comprehensive derivation appears in Appendix~\ref{app:derivation2}. Although Eq.~\ref{eq:objective} invokes the Bradley-Terry model, parallel results extend to the broader class of Plackett-Luce models~\citep{plackett1975analysis, luce2012individual}, as developed in Appendix~\ref{app:plackett_luce_models}.

With the likelihood of observed human preferences now expressed in terms of the optimal policy rather than a reward estimator, we are positioned to formulate a maximum likelihood training objective for a parameterized policy $\pi_\theta$. In analogy to the reward modeling strategy (see Eq.~\ref{eq:reward_model}), this yields the policy learning objective:

\begin{equation}\label{eq:optimum_model}
    \mathcal{L}_\text{DPO}(\pi_{\theta}; \pi_\text{ref}) = -\mathbb{E}_{(x, y_w, y_l)\sim \mathcal{D}}\left[\log \sigma \left(\beta \log \frac{\pi_{\theta}(y_w

In this formulation, we estimate an implicit reward through an alternative parameterization, where the optimal policy corresponds directly to $\pi_\theta$. Additionally, because our approach mirrors fitting a reparameterized Bradley-Terry model, it inherits certain theoretical guarantees, such as consistency under appropriate assumptions about the underlying preference data distribution \cite{bong2022generalized}. Further analysis of DPO’s theoretical characteristics in comparison to related methods is provided in Section~\ref{sec:theory}.

\textbf{Mechanistic analysis of the DPO update.} To gain insight into the operational dynamics of DPO, one can study the gradient of the associated loss, $\mathcal{L}_\text{DPO}$. With respect to the parameters $\theta$, the gradient is expressed as:

\begin{multline*}\label{eq:gradient}
    \nabla_\theta \mathcal{L}_\text{DPO}(\pi_\theta;\pi_\text{ref}) = \\ -\beta\mathbb{E}_{(x, y_w, y_l) \sim \mathcal{D}} \bigg[\underbrace{\sigma(\hat{r}_\theta(x, y_l) - \hat{r}_\theta (x, y_w))}_\text{greater magnitude when reward ranking is incorrect}\bigg[\underbrace{\nabla_\theta\log \pi(y_w \mid x)}_\text{push up $y_w$ likelihood} - \underbrace{\nabla_\theta\log\pi(y_l \mid x)}_\text{push down $y_l$ likelihood}\bigg]\bigg],
\end{multline*}

where the implicit reward function is given by $\hat{r}_\theta(x, y) = \beta \log \frac{\pi_\theta(y \mid x)}{\pi_\text{ref}(y \mid x)}$, with further elaboration in Section~\ref{sec:theory}. Intuitively, the gradient of $\mathcal{L}_\text{DPO}$ works by encouraging an increase in the probability assigned to favored completions $y_w$, while discouraging the probability for less desirable completions $y_l$. A key feature is that each data point is weighted by the extent to which the implicit reward model $\hat{r}_\theta$ overvalues the disfavored completion, multiplied by $\beta$—that is, the degree of misranking produced by the reward model, modulated by the KL regularization strength. Empirically, we observe that this form of weighting is crucial: eliminating the weighting factor (using a naïve variant of the update) leads to language model collapse, as evidenced in Appendix Table~\ref{tab:unlikelihood_generations}.

\textbf{DPO overview.} The typical DPO procedure proceeds as follows: 1) For each prompt $x$, generate candidate completions $y_1, y_2$ using $\pi_\text{ref}(\cdot \mid x)$, then annotate these outputs according to human preferences to form an offline preference dataset $\mathcal{D} = \{x^{(i)}, y_w^{(i)}, y_l^{(i)}\}_{i=1}^N$; 2) update the language model $\pi_\theta$ by minimizing the DPO loss $\mathcal{L}_\text{DPO}$, conditioned on the specified reference policy $\pi_\text{ref}$, dataset $\mathcal{D}$, and target $\beta$ parameter. 
In realistic scenarios, rather than generating new samples and annotations, one often relies on pre-existing public preference datasets. Because such datasets have completions generated by $\pi^\text{SFT}$, we typically set $\pi_\text{ref} = \pi^\text{SFT}$ whenever this is feasible. If $\pi^\text{SFT}$ is not accessible, $\pi_\text{ref}$ is initialized instead by fitting it to maximize the likelihood of preferred completions ${(x, y_w)}$, that is, ${\pi_\text{ref} = \argmax_{\pi}\mathbb{E}_{x, y_w \sim \mathcal{D}}\left[\log \pi(y_w \mid x)\right]}$. This step is designed to reduce the mismatch between the ideal, but unavailable, reference distribution and the practical $\pi_\text{ref}$ utilized during DPO. Implementation specifics and hyperparameter choices are further elaborated in Appendix~\ref{app:implementation}.

\section{Theoretical Analysis of DPO} This section aims to deepen the interpretation of the DPO framework, offer supporting theoretical results, and clarify how DPO overcomes certain limitations of actor-critic methods employed in RLHF, such as PPO~\cite{schulman2017proximal}.

\label{sec:theory}

\subsection{Your Language Model Is Secretly a Reward Model} DPO sidesteps the need for an explicit reward function and a separate RL phase for policy learning, instead utilizing a unified maximum likelihood loss. Observe that the objective in Eq.~\ref{eq:main_eq} is mathematically equivalent to applying a Bradley-Terry model, where the reward function is parameterized as $r^*(x, y) = \beta \log\frac{\pi^*_\theta(y \mid x)}{\pi_\text{ref}(y \mid x)}$. The parameterized model $\pi_\theta$ is then optimized, paralleling the optimization used for reward models as in Eq.~\ref{eq:reward_model}, after a suitable variable substitution. This section will develop the theoretical underpinnings of this parameterization, demonstrate that it imposes no restrictions on the family of recoverable reward models, and ensures that the optimal policy can be exactly obtained. We start by formalizing an equivalence relation between reward functions.

\begin{definition}
We define two reward functions $r(x, y)$ and $r'(x, y)$ as equivalent if and only if there exists some function $f$ such that ${r(x, y)-r'(x, y) = f(x)}$.     
\end{definition}

This definition clearly establishes an equivalence relation, thus dividing the collection of reward functions into equivalence classes. We can formalize the following results:

\begin{lemma}\label{lemma:same_prefrence} Within the Plackett-Luce, and specifically the Bradley-Terry, preference frameworks, any two reward functions belonging to the same equivalence class will generate identical preference distributions.
\end{lemma}

\begin{lemma}\label{lemma:same_policy}
    Any pair of reward functions within a shared equivalence class will induce the same optimal policy in the corresponding constrained reinforcement learning problem.
\end{lemma}

The proofs for these lemmas are straightforward and can be found in Appendix \ref{app:lemma1}. The first lemma highlights a well-documented issue of under-identification inherent in models from the Plackett-Luce family \cite{plackett1975analysis}. Owing to this lack of identifiability, it is common practice to enforce additional constraints to ensure that maximum likelihood estimators, as described in Eq. \ref{eq:reward_model}, are well-defined \cite{bong2022generalized}. The second lemma implies that all reward functions contained within a single equivalence class will determine the same optimal policy. Consequently, our ultimate objective is to recover any representative reward function from the optimal equivalence class. The following theorem, whose proof is given in Appendix~\ref{app:thm1}, formalizes this:

\begin{theorem}\label{thm:main}
    Provided certain regularity conditions hold, each reward equivalence class consistent with the Plackett-Luce (including the Bradley-Terry model) admits a representation of the form ${r(x, y) = \beta \log \frac{\pi(y\mid x)}{\pi_\text{ref}(y\mid x)}}$ for some model $\pi(y\mid x)$ and a given reference model $\pi_\text{ref}(y \mid x)$.
\end{theorem}

\begin{sproof}
    Take any reward function $r(x, y)$ which gives rise to an optimal policy $\pi_r(y \mid x)$, as defined by Eq. \ref{eq:op_policy}. We aim to demonstrate that within the equivalence class of $r$, there exists a reward function expressible through the stated reparameterization. Consider the mapping $f$ defined by  
\begin{equation}
    f(r; \pi_\text{ref}, \beta)(x, y) = r(x, y) - \beta\log\sum_{y}\pi_\text{ref}(y\mid x)\exp\left(\frac{1}{\beta}r(x, y)\right)
\end{equation}
Here, $f$ operates by normalizing the reward function via the log-partition function of $\pi_r$. As the normalization term depends solely on $x$, $f(r; \pi_\text{ref}, \beta)(x, y)$ is guaranteed to remain within the same equivalence class as $r(x, y)$. Substituting $r$ by the right-hand side of Eq.~\ref{eq:main_eq} (which is universally valid for any reward function) yields $f(r; \pi_\text{ref}, \beta)(x, y) = \beta \log \frac{\pi_r(y\mid x)}{\pi_\text{ref}(y\mid x)}$. Therefore, $f$ constructs a representative of the equivalence class of $r$ possessing the required structure, confirming that our proposed reparameterization imposes no restrictions on the reward model.
\end{sproof}

An alternative perspective on Theorem~\ref{thm:main} is that it precisely determines which reward function is chosen within each equivalence class by the DPO reparameterization—namely, the one that fulfills:

\begin{equation}\label{eq:lag_p}
     \sum_{y}\underbrace{\pi_\text{ref}(y\mid x)\exp\left(\frac{1}{\beta}r(x, y)\right)}_{=\pi(y\mid x)\text{, using Thm.~\ref{thm:main} reparam.}} = 1,
\end{equation}

This requirement guarantees that $\pi(y\mid x)$ constitutes a legitimate probability distribution (with all probabilities non-negative and summing to one). Importantly, when relating this to Eq.~\ref{eq:op_policy}, we recognize that Eq.~\ref{eq:lag_p} defines the partition function associated with the optimal policy that corresponds to the reward $r(x, y)$. The central insight behind the DPO method is to enforce particular constraints on the otherwise underdetermined Plackett-Luce class of preference models (including the Bradley-Terry variant). These constraints do not limit the expressivity of reward models but, crucially, they allow the optimal policy described by Eq.~\ref{eq:op_policy} to be expressed in closed form for any prompt $x$.

\subsection{Instability of Actor-Critic Algorithms}
Our framework also facilitates the identification of instabilities in commonly used actor-critic methods for RLHF, such as PPO. Consistent with the RLHF workflow, we analyze the RL fine-tuning stage discussed in Section~\ref{section:prelims}. This connects to the control-as-inference formalism \cite{levine2018reinforcement} for the constrained RL formulation given by \ref{eq:RL}. Considering a parameterized policy $\pi_{\theta}(y\mid x)$, the typical objective is to minimize the divergence $\mathbb{D}_{\text{KL}}[\pi_{\theta}(y|x) \mid \mid \pi^*(y\mid x)]$, where $\pi^*$ represents the optimal policy from Eq.~\ref{eq:optimum_model} derived from the reward $r_{\phi}(y, x)$. Upon simplification, this results in the following optimization problem:

\begin{equation}\label{eq:AC}
    \max_{\pi_{\theta}}\mathbb{E}_{\pi_{\theta}(y\mid x)}\left[\underbrace{r_{\phi}(x, y) -\beta\log\sum_{y}\pi_\text{ref}(y\mid x)\exp\left(\frac{1}{\beta}r_{\phi}(x, y)\right)}_{f(r_{\phi}, \pi_\text{ref}, \beta)} - \underbrace{\beta\log\frac{\pi_{\theta}(y\mid x)}{\pi_\text{ref}(y\mid x)}}_{\text{KL}}\right]
\end{equation}

This objective has been the focus of optimization in previous research 
\citep{ziegler2020finetuning, stiennon2022learning, bai2022training, ouyang2022training}, where the DPO-equivalent reward is used for the reward class $r_{\phi}$. Within this framework, the normalization component of $f(r_{\phi}, \pi_\text{ref}, \beta)$ can be viewed as the soft value function corresponding to the reference policy $\pi_\text{ref}$. Although this term does not alter the optimality of the solution, omitting it can lead to high-variance policy gradient estimates, thus impeding stable learning. One strategy to address this normalization is to employ a learned value function; however, optimizing such a function can pose its own challenges. Prior approaches have sometimes replaced the normalization term with a reward normalized by a human-generated completion, effectively serving as a single-sample Monte Carlo approximation. Distinctively, the DPO reparameterization leads to a reward function formulation that eliminates the need for any explicit baselines. 

\section{Experiments}
This section presents an empirical assessment of DPO’s capability to train policies using preference data directly. We begin with controlled experiments in text generation, investigating DPO’s efficiency in balancing reward maximization against KL-divergence minimization with respect to the reference policy, and compare its performance to standard preference optimization algorithms like PPO. Subsequently, we extend the evaluation to larger models and more challenging RLHF scenarios, such as dialogue and summarization tasks. Across these settings, we observe that DPO often matches or surpasses strong baselines like RLHF with PPO or the strategy of selecting the best out of $N$ sampled trajectories based on a learned reward, all with minimal hyperparameter tuning. Preceding the presentation of our empirical findings, we outline the experimental methodology; supplementary details can be found in Appendix~\ref{app:exp_details}.

\textbf{Tasks.} Our empirical evaluation covers three distinct open-ended text generation tasks. Across all experiments, methods are trained to learn a policy from a preference dataset denoted by $\mathcal{D}=\bigl\{x^{(i)}, y_w^{(i)}, y_l^{(i)}\bigr\}_{i=1}^N$. In the case of \textbf{controlled sentiment generation}, $x$ corresponds to a movie review prefix sourced from the IMDb dataset \cite{maas-EtAl:2011:ACL-HLT2011}, with the objective that the policy generates a continuation $y$ that exhibits positive sentiment. For rigorous assessment, we synthetically construct preference pairs for generations by leveraging a pre-trained sentiment classifier, selecting pairs such that $p(\text{positive}\mid x,y_w)>p(\text{positive}\mid x,y_l)$. When training SFT models, GPT-2-large is fine-tuned to convergence using the IMDB training split (detailed in App~\ref{app:sentiment_details}). In the \textbf{summarization} task, $x$ is a Reddit forum post, and the policy is expected to produce a concise summary $y$ encapsulating the post’s main ideas. Consistent with established benchmarks, the Reddit TL;DR summarization corpus \citep{volske-etal-2017-tl} is employed along with preference annotations collected by \citeauthor{stiennon2022learning}. For RLHF, we utilize an SFT model specifically adapted on human-authored summaries of forum posts,\footnote{\url{https://huggingface.co/CarperAI/openai_summarize_tldr_sft}} making use of the TRLX framework \citep{leandro_von_werra_2023_7790115}. Note that the human preference labels originate from a separate, but equivalently trained, SFT model as described in \citeauthor{stiennon2022learning}. Lastly, the \textbf{single-turn dialogue} task considers $x$ as a user input, which could range from a scientific inquiry to advice-seeking. The objective is for the policy to craft a response $y$ that is both informative and engaging, drawing upon the Anthropic Helpful and Harmless dialogue dataset \citep{bai2022training} comprising 170,000 interactions between humans and automated agents. Each record concludes with two candidate completions from a large (identity unspecified) language model, annotated to indicate human preference. As a pre-trained SFT model is not provided in this domain, we develop an SFT baseline by fine-tuning a publicly available language model exclusively on preferred completions.

\textbf{Evaluation.} Our experimental framework employs two complementary evaluation strategies. For the controlled sentiment generation scenario, we assess algorithmic performance in constrained reward maximization by measuring the tradeoff between achieved reward and KL-divergence from the reference policy. This analysis is enabled by access to a ground-truth reward signal, specifically a sentiment classifier. Conversely, in practical, real-world deployments where the true reward function remains unknown, we adopt a \textit{win rate} metric relative to a baseline policy. In these settings—summarization and single-turn dialogue—we rely on GPT-4 as a stand-in for human raters, using it to judge the quality of summaries and the helpfulness of responses. The baselines consist of reference summaries from the test split for summarization and the human-preferred responses in the dialogue test set. Though recent literature suggests that large language models may surpass conventional metrics as automatic evaluators \citep{Chen2023ExploringTU}, we bolster our choice of GPT-4 with a human annotation study, as detailed in Sec.~\ref{sec:human-judgments}. Our results indicate a strong alignment between GPT-4 and human judgments, with levels of agreement often matching or exceeding those found between individual human annotators.

\textbf{Methods.} Alongside DPO, we examine a suite of established methods for aligning language models with human preferences. For the summarization domain, we test zero-shot prompting with \textbf{GPT-J} \citep{gpt-j}, while in dialogue we utilize 2-shot prompting with \textbf{Pythia-2.8B} \citep{biderman2023pythia}. Furthermore, we assess both the \textbf{SFT} model and a variant called \textbf{Preferred-FT}, which is fine-tuned via supervised learning using selected completions $y_w$—sourced from SFT outputs (in sentiment and summarization) or from a generic language model (in dialogue). Another pseudo-supervised alternative is \textbf{Unlikelihood}~\citep{welleck2019neural}, where optimization involves promoting the probability of $y_w$ and suppressing the probability of $y_l$, weighted by an optional coefficient $\alpha\in[0,1]$ applied to the unlikelihood loss component. We also implement \textbf{PPO} \citep{schulman2017proximal} with reward signals inferred from preference data, in addition to \textbf{PPO-GT}—an oracle variant leveraging the ground-truth reward available in the sentiment setting. For the latter, two versions are considered: an out-of-the-box implementation \cite{leandro_von_werra_2023_7790115} and a customized approach with reward normalization and optimized hyperparameters for enhanced outcomes (similar modifications are applied when running PPO with learned rewards). Lastly, the \textbf{Best of $N$} baseline is included: here, $N$ samples are generated from the SFT model (or Preferred-FT in the dialogue scenario), and the candidate scoring highest under a learned reward model is selected. While this method decouples reward model quality from PPO optimization and generally yields strong performance, it is computationally prohibitive at test time for even modest values of $N$, due to the need for multiple samplings per prompt.

\subsection{How well can DPO optimize the RLHF objective?}

In standard RLHF methods, the objective of maximizing reward subject to a KL constraint serves to encourage reward maximization while simultaneously limiting policy divergence from a given reference policy. Accordingly, a fair comparison between different algorithms should jointly consider both attained reward and the corresponding KL divergence; attaining marginally greater reward at the cost of a substantially increased KL is often undesirable. Figure~\ref{fig:frontier-tldr-main} illustrates the tradeoff curve between reward and KL for a range of algorithms within the sentiment task. For each method, we conduct several independent training runs, each corresponding to a distinct policy conservativeness hyperparameter (for PPO: target KL $\in\{3,6,9,12\}$; for DPO: $\beta \in \{0.05,0.1,1,5\}$; for unlikelihood: $\alpha\in\{0.05,0.1,0.5,1\}$; preferred-FT: different random seeds), amounting to a total of 22 trials. At every 100 training steps up to convergence, we assess the policies using a fixed set of test prompts, calculating both the mean reward according to the true reward metric and the average sequence-level KL\footnote{Specifically, this refers to the aggregate of timestep-wise KL-divergences.} with respect to the reference policy, $\text{KL}\left(\pi\mid \mid \pi_\text{ref}\right)$. The results show that DPO consistently yields the most favorable tradeoff curve, securing maximal reward with lower KL values relative to other baselines. This finding is significant for several reasons. Firstly, both DPO and PPO are designed to optimize an identical objective, yet DPO does so with markedly greater efficiency; DPO strictly surpasses PPO in terms of the reward versus KL balance. Secondly, DPO attains a superior tradeoff curve not only in comparison to standard PPO but also in the case where PPO has access to the true reward signal (PPO-GT).

\subsection{Can DPO scale to real preference datasets?}
\label{sec:dpo-real-datasets}
We proceed to assess the effectiveness of DPO fine-tuning using real-world preference data in tasks such as summarization and single-turn dialogue. In the domain of summarization, conventional automated metrics like ROUGE often show weak alignment with human judgments~\citep{stiennon2022learning}. Previous research has indicated that applying PPO for fine-tuning language models based on human feedback can yield higher-quality summaries. For our experiments, we evaluate different fine-tuning strategies by generating outputs on the test split of the TL;DR summarization dataset and calculating the mean win rate when compared with the reference summaries in the test partition. For each approach, we sample completions at temperatures ranging from 0.0 to 1.0; the resulting win rates are illustrated in Figure~\ref{fig:frontier-tldr-main} (right panel). DPO, PPO, and Preferred-FT methods all start from the same GPT-J SFT model\footnote{\url{https://huggingface.co/CarperAI/openai_summarize_tldr_sft}} for fine-tuning. Our results indicate that DPO achieves a win rate of roughly 61\% at a sampling temperature of 0.0, surpassing PPO, which reaches about 57\% under its optimal temperature conditions. Additionally, DPO attains a greater peak win rate in comparison to the best of $N$ baseline. It is important to mention that we did not extensively optimize the $\beta$ hyperparameter for DPO; thus, these outcomes could be a conservative estimate of DPO's capabilities. Furthermore, DPO demonstrates considerably greater resilience to variations in sampling temperature relative to PPO, whose performance can deteriorate to that of the original GPT-J model as temperature increases. Preferred-FT does not produce notable gains over the SFT baseline. We further directly compare DPO and PPO in human preference studies in Section~\ref{sec:human-judgments}, showing that at a temperature of 0.25, DPO samples are selected over PPO samples 58\% of the time.

For single-turn dialogue tasks, we conduct evaluations on a segment of the Anthropic HH dataset’s test partition \citep{bai2022training}, focusing specifically on interactions involving only one exchange between a human and an assistant. In the case of GPT-4, we utilize the test set’s preferred completions as the ground truth to determine win rates across methods. Since a standardized SFT model tailored to this dataset does not exist, we initialize our experiments using the pre-trained Pythia-2.8B model. From there, we employ Preferred-FT to fine-tune a reference model on selected completions, ensuring that the outputs remain representative of the underlying data distribution, before proceeding to DPO training. 

Additionally, our experiments include comparisons with the Best of 128 Preferred-FT completions (established as the optimal plateau point for this baseline; see Appendix Figure~\ref{fig:best-of-n}) and a 2-shot prompting setup using the base Pythia-2.8B model. Across a range of optimal sampling temperatures, DPO demonstrates equivalent or superior results compared to the other methods. Furthermore, we assess an RLHF model trained via PPO on the Anthropic HH dataset \footnote{\url{https://huggingface.co/reciprocate/ppo_hh_pythia-6B}}, sourced from a widely used implementation \footnote{\url{https://github.com/CarperAI/trlx/tree/main/examples/hh}}. However, we are unable to identify any prompting or temperature configuration that achieves better performance than the base Pythia-2.8B. 

Given our findings with the TL;DR benchmark and the observation that both PPO and Best of 128 aim to optimize an identical reward signal, we regard Best of 128 as an approximate stand-in for PPO-level results. Collectively, our study finds DPO to be the only approach that, while remaining computationally efficient, reliably outperforms the preferred completions within the Anthropic HH dataset. Furthermore, it attains comparable or superior results relative to the more computationally intensive Best of 128. As illustrated in Figure~\ref{fig:dialogue-main}, DPO also achieves its peak effectiveness after relatively few iterations.

\subsection{Generalization to a new input distribution}

To further assess how PPO and DPO respond to distributional changes, we test the policies trained on Reddit TL;DR summarization tasks against an out-of-domain dataset: the test split of CNN/DailyMail news articles \citep{nallapati-etal-2016-abstractive}. For these evaluations, we employ the optimal sampling temperatures previously identified for TL;DR (0 and 0.25). Results are detailed in Table~\ref{tab:ood}. We measured GPT-4’s win rates using the same prompt applied for Reddit TL;DR, with the phrase ``forum post'' adapted to ``news article'' as appropriate. On this novel input distribution, DPO consistently yields substantially higher win rates than PPO. These findings offer preliminary support that DPO-trained policies retain generalization capabilities comparable to those of PPO, despite DPO lacking access to the supplementary unlabeled Reddit TL;DR prompts utilized by PPO.

\subsection{Validating GPT-4 judgments with human judgments}
\label{sec:human-judgments}
We carry out a human evaluation to establish the consistency of GPT-4’s preference judgments, focusing on outputs from the TL;DR summarization experiment and utilizing two distinct GPT-4 prompts. The \textbf{GPT-4 (S)} (simple) prompt only requests identification of which summary better captures the core information in the post, while the \textbf{GPT-4 (C)} (concise) prompt further asks which summary is more concise. We include the latter since we observe that with the \textbf{GPT-4 (S)} prompt, GPT-4 often favors lengthier and more repetitive summaries than those preferred by human raters. Appendix~\ref{app:prompts} provides the full prompt texts. To sample a broad range of summary qualities, we select the best-performing (DPO, temp. 0.25), the least effective (PPO, temp. 1.0), and an intermediate (SFT, temp. 0.25) method for pairwise comparisons against the PPO model sampled greedily (using its optimal temperature). Across both prompt variants, GPT-4's choices correlate with human selections at rates comparable to human-human agreement, indicating that GPT-4 serves as a viable stand-in for human evaluation (note that due to limited human participants, multiple judgments are collected only for the DPO and PPO-1 matchups). Notably, the \textbf{GPT-4 (C)} prompt generates win rates more closely aligned with human preferences, so we adopt this prompt for reporting the main results in Section~\ref{sec:dpo-real-datasets}. Further methodological specifics, such as the rater interface and a roster of participants, can be found in Appendix~\ref{app:human-study}.

\section{Discussion}
Preference-based learning offers an effective and scalable approach for developing language models that are both capable and aligned with user intentions. In this work, we have presented DPO, a straightforward training framework that enables language models to learn from preferences without employing reinforcement learning. Unlike traditional methods that recast preference learning as a reinforcement learning task in order to utilize established RL techniques, DPO instead establishes a direct correspondence between language model policies and reward functions. This allows language models to align with human preferences using a basic cross-entropy loss, thereby avoiding reinforcement learning altogether and maintaining full generality. Remarkably, DPO achieves comparable or superior performance to existing RLHF approaches, including those utilizing PPO, with minimal hyperparameter adjustment; this considerably lowers the effort required to train language models based on human preferences.

\textbf{Limitations \& Future Directions.} Our findings open up several avenues for future research. One pertinent question concerns the out-of-distribution generalization of the DPO policy as opposed to approaches relying on explicit reward functions. Preliminary evidence indicates that DPO-trained policies generalize on par with PPO-based counterparts, though more extensive investigations are warranted. For instance, it is worth exploring whether employing self-labeling mechanisms with the DPO policy allows for efficient utilization of prompts without preference labels. Another line of inquiry involves understanding how reward over-optimization may appear within direct preference optimization, and whether the observed minor dip in performance (Figure~\ref{fig:dialogue-main}-right) is a consequence of such effects. Furthermore, although our assessments include models with up to 6B parameters, scaling DPO to much larger, cutting-edge models remains an intriguing and promising research direction. In the context of evaluation, we observe that GPT-4-derived win rates are sensitive to the phrasing of prompts; further work could refine strategies for obtaining reliable automated evaluations. Lastly, the versatility of DPO suggests its potential for application in training generative models beyond language, spanning additional data modalities.